\section{Conclusion and Future Work}

\subsection{Achievements}

This research has successfully accomplished all primary objectives through the development and validation of a novel Multi-layer Probability Fusion Algorithm based on relative percentile ranking for football transfer prediction. The algorithmic development phase resulted in a comprehensive percentile-based feature engineering system that transforms absolute performance metrics into contextually-aware relative rankings, successfully addressing the fundamental challenge of cross-context player comparison in sports analytics. The implementation of a multi-algorithm optimisation framework incorporating Genetic Algorithm, Particle Swarm Optimisation, Simulated Annealing, and Random Search demonstrated the effectiveness of systematic parameter tuning approaches, achieving optimal weight configurations through rigorous comparative analysis.

The empirical validation demonstrated exceptional performance with the PSO algorithm achieving 92.8\% accuracy with 96.48\% PR-AUC on the optimization framework, whilst maintaining 77.8\% prediction accuracy on Inter Milan's 2023-2024 transfer outcomes, successfully predicting 21 out of 27 player decisions through the optimised algorithm. The research established the superiority of relative percentile ranking over traditional absolute metric approaches through comprehensive statistical significance testing, with all optimisation algorithms showing statistically significant improvements (p < 0.001) over baseline methods. The systematic validation framework incorporating temporal cross-validation, 5-run multi-run experiments, and ground truth verification ensures the robustness and reliability of the research findings.

\subsection{Innovations}

The core innovation of this research lies in the development of a comprehensive percentile-based feature engineering approach that represents a fundamental paradigm shift from absolute performance evaluation to contextually-aware relative assessment. This methodology addresses critical limitations in existing sports analytics by eliminating systematic biases introduced by varying competitive contexts, enabling fair cross-player comparisons within league-wide performance distributions, and creating position-specific evaluation criteria that accurately reflect tactical role requirements.

The implementation of an enhanced multi-algorithm weight optimisation framework represents a significant methodological advancement through the integration of four complementary optimisation approaches. The systematic comparison of bio-inspired algorithms (Genetic Algorithm, Particle Swarm Optimisation) with traditional approaches (Simulated Annealing, Random Search) establishes best practices for integrated parameter optimisation in sports prediction problems, with the modular design ensuring extensibility for future research applications.

The development of a composite evaluation metric combining PR-AUC, F1-Score, Balanced Accuracy, and Brier Score addresses the fundamental challenge of multi-objective optimisation in sports prediction contexts. This innovation balances ranking capability with classification performance whilst accounting for class imbalance common in transfer prediction problems and incorporating probability calibration quality essential for practical decision-making. The composite metric design provides comprehensive performance assessment beyond simple accuracy measures, enabling optimisation processes that excel across different evaluation dimensions rather than maximising single metrics that may not reflect real-world decision requirements.

\subsection{Limitations}

The main limitations of this study are as follows: The research focus on Inter Milan as a single club case study, whilst providing detailed analysis of transfer decision patterns, inherently limits the broader applicability of the algorithm across different tactical systems, financial contexts, and club cultures. The limited sample size of 25 players constrains statistical power for certain analyses and may not capture the full spectrum of player departure scenarios encountered across different club contexts. The temporal constraints of single-season optimisation may fail to capture longer-term performance patterns and market dynamics that influence transfer decisions, whilst the specific economic conditions and regulatory environment during the study period may not represent typical transfer market conditions.

\subsection{Future Work}

Future research directions encompass extensions to validate and enhance the algorithm's applicability across broader contexts. The most critical immediate priority involves multi-club validation across different Serie A clubs to establish league-specific applicability, followed by cross-league validation incorporating Premier League, La Liga, and Bundesliga data to assess algorithm performance across varying competitive environments and tactical approaches. Temporal extension through multi-season validation would strengthen the research foundation by assessing algorithm stability over time and investigating the integration of performance trend analysis for improved prediction accuracy.

Algorithmic enhancements present significant opportunities for improved performance through advanced machine learning integration, including deep learning architectures for complex pattern recognition and ensemble methods that combine multiple prediction approaches. The incorporation of multi-modal data sources represents a particularly promising research direction, for example,social media sentiment analysis for player satisfaction indicators. Natural language processing integration could enable the algorithm to incorporate news sentiment and social media patterns that may provide early indicators of transfer intentions.

The algorithmic framework demonstrates potential for broader applications beyond football, including basketball player movement and trade predictions, cricket player selection and performance optimisation, and general sports talent identification and development.